\documentclass[11pt,a4paper]{article}

% Viane chronic-pain experience-sampling manuscript 02: the flow construct.
% Compile with: tectonic main.tex
% References use natbib with the bundled apalike BibTeX style, so no biber dependency is required.
\PassOptionsToPackage{colorlinks=true,linkcolor=NavyBlue,citecolor=NavyBlue,urlcolor=NavyBlue}{hyperref}
\usepackage[a4paper,margin=2.54cm,top=2.85cm,bottom=2.7cm,headheight=14pt,headsep=16pt]{geometry}
\usepackage[T1]{fontenc}
\usepackage[utf8]{inputenc}
\usepackage{newtxtext,newtxmath}
\usepackage{microtype}
\usepackage{setspace}
\usepackage{booktabs}
\usepackage{array}
\usepackage{tabularx}
\usepackage{graphicx}
\usepackage{caption}
\usepackage{float}
\usepackage{fancyhdr}
\usepackage{amsmath}
\usepackage{enumitem}
\usepackage[dvipsnames]{xcolor}
\usepackage{xurl}
\usepackage[american]{babel}
\usepackage[round]{natbib}
\usepackage{orcidlink}
\usepackage{hyperref}
\usepackage{cleveref}

\AtBeginDocument{%
  \providecommand{\doi}[1]{\url{https://doi.org/#1}}%
  \renewcommand{\doi}[1]{\url{https://doi.org/#1}}%
}
\captionsetup{font=small,labelfont=bf,labelsep=period,justification=raggedright,singlelinecheck=false}
\setlength{\parindent}{1.5em}
\setlength{\textfloatsep}{14pt plus 3pt minus 3pt}
\setlength{\floatsep}{12pt plus 2pt minus 2pt}
\setlength{\intextsep}{12pt plus 2pt minus 2pt}
\onehalfspacing
\pagestyle{fancy}
\fancyhf{}
\fancyhead[L]{\footnotesize Decomposing Flow in Daily Life}
\fancyhead[R]{\footnotesize Van Severen, Viane, De Paepe, \& Crombez}
\fancyfoot[C]{\footnotesize\thepage}
\renewcommand{\headrulewidth}{0.4pt}
\renewcommand{\footrulewidth}{0pt}
\fancypagestyle{plain}{%
  \fancyhf{}%
  \fancyhead[L]{\footnotesize Decomposing Flow in Daily Life}%
  \fancyhead[R]{\footnotesize Van Severen, Viane, De Paepe, \& Crombez}%
  \fancyfoot[C]{\footnotesize\thepage}%
  \renewcommand{\headrulewidth}{0.4pt}%
}
\makeatletter
\setlength{\@fptop}{0pt}
\setlength{\@fpsep}{12pt plus 0fil}
\setlength{\@fpbot}{0pt plus 1fil}
\renewcommand{\topfraction}{0.94}
\renewcommand{\bottomfraction}{0.80}
\renewcommand{\textfraction}{0.06}
\renewcommand{\floatpagefraction}{0.82}
\setcounter{topnumber}{3}
\setcounter{bottomnumber}{2}
\setcounter{totalnumber}{4}
\makeatother

\newcommand{\figurenote}[1]{\par\vspace{2pt}\footnotesize\textit{Note.} #1}
\newcommand{\asset}[1]{../../assets/02_manuscript_FlowInducedAnalgesia/#1}
\title{Flow Is Not More Than Its Parts: Decomposing the Challenge-Skill Model and Its Momentary Association with Chronic Pain}
\author{Stijn Van Severen*, Ilse Viane, Annick De Paepe, and Geert Crombez}
\date{}

\begin{document}

\begin{titlepage}
\thispagestyle{fancy}
\vspace*{\fill}
\begin{center}
{\fontsize{18.2}{22.8}\selectfont\textbf{Flow Is Not More Than Its Parts:\\
Decomposing the Challenge-Skill Model and Its\\
Momentary Association with Chronic Pain}\par}
\vspace{1.2em}
\normalsize
Stijn Van Severen\textsuperscript{1,*}\,\orcidlink{0009-0003-9181-966X},
Ilse Viane\textsuperscript{1},
Annick De Paepe\textsuperscript{1}\,\orcidlink{0000-0002-9366-3692}, and
Geert Crombez\textsuperscript{1}\,\orcidlink{0000-0002-4744-8561}\\[0.6em]
\footnotesize
\textsuperscript{1}\,Department of Experimental-Clinical and Health Psychology,
Ghent University, Ghent, Belgium\\
\textsuperscript{*}\,Corresponding author: \href{mailto:stijn.vanseveren@ugent.be}{stijn.vanseveren@ugent.be}
\end{center}
\vspace{0.8em}

\begin{abstract}
\noindent
Flow, the absorbed and intrinsically rewarding state said to arise when perceived challenge and
perceived skill are both high and matched, has been proposed as a route to analgesia. The
proposal rests on two structural assumptions that are built into experience-sampling measures of
flow but are rarely tested: (1) that the challenge-skill condition acts through the \emph{match}
between its two components, and (2) that the constituent appraisals form one construct that a
composite summarizes fairly. We tested both assumptions, and the momentary association with pain, in
$68$ patients with chronic pain who rated challenge, skill, absorption, and enjoyment together
with pain, its interference, attention to pain, and its threat value eight times per day for two
weeks ($5{,}976$ moments with all four appraisals, median $=92$ per person). Neither assumption
held. Balance, the negative absolute challenge-skill discrepancy, added nothing to an additive
model of challenge and skill ($\chi^2(1)=0.56$, $p=.45$), and neither did a one-sided kink at
equality, a multiplicative interaction, or a full second-order response surface (all $p \geq .27$);
between persons, better-matched activities went with \emph{less} flow experience ($b=-.111$,
$p=.028$). Entering the four constituents jointly rather than as a composite improved fit for
every pain outcome (all $p<.001$), because absorption and skill tracked lower pain
($b=-.055$ and $b=-.045$) while challenge reversed sign and tracked \emph{higher} pain
($b=+.032$, $p=.004$). The flow experience was concurrently associated with lower pain
interference ($b=-.139$), pain ($b=-.094$), threat ($b=-.079$), and attention to pain
($b=-.074$, all $p<.001$), but no lag-1 effect reached significance in either direction while
autoregressive effects remained large, and the least painful region of the challenge-skill plane
was the low-challenge, high-skill relaxation channel rather than flow. Flow in daily life with
chronic pain is best described by how absorbing and how manageable an activity is, not by how
well demand and capacity are matched; the composite adds a label where the components carry the
information.\\[0.5em]
\noindent\textit{Keywords:} flow; challenge-skill balance; chronic pain; constructed measures;
ecological momentary assessment; multilevel modelling
\end{abstract}

\vspace*{\fill}
\end{titlepage}
\setcounter{page}{2}
\section{Introduction}

\subsection{Flow, and the proposal that it relieves pain}
Flow is the state of complete absorption in an activity that is experienced as intrinsically
rewarding and that arises, on the classic account, when the demands of the activity and the
person's capacity to meet them are both high and in balance
\citep{csikszentmihalyi1990flow,nakamura2014flow}. The construct was built for
experience sampling from the start: challenge and skill were rated at signalled moments, and the
resulting challenge-skill plane was divided into four channels, flow, anxiety, relaxation, and
apathy, whose quality of experience was then compared
\citep{massimini1988systematic,carli1988quality,csikszentmihalyi1987validity}.

Interest in flow has recently reached pain research. \citet{gromakovskis2025flow} review evidence
for flow-induced analgesia, the proposal that deep absorption in a competing activity blunts pain,
and place it alongside the biopsychosocial and neurobiological accounts of pain modulation. The
proposal is attractive because it converges with an established idea in pain psychology. In the
cognitive-affective model of the interruptive function of pain \citep{eccleston1999pain}, pain and
current goals compete for a limited attentional capacity, and the outcome of that competition
depends on the threat value of the pain and on the characteristics of the activity it interrupts
\citep{crombez1998catastrophizing,crombez1999fear,vandamme2010keeping}. Threat-related appraisal and
monitoring amplify the attentional pull of pain
\citep{crombez2013attention,eccleston2007worry}, and the same appraisals sustain avoidance of
demanding activity \citep{crombez2012fear}. Flow describes the opposite pole of the
same competition: an activity so absorbing that competing signals fail to intrude. If that is
right, flow states should be states of reduced pain.

The evidence adduced for flow-induced analgesia comes almost entirely from paradigms in which
flow is induced, through virtual reality, music, or structured tasks, and pain is measured during
the induction \citep{gromakovskis2025flow}. Such designs support causal inference but say little
about whether the same relation holds in the unstructured daily life of patients whose pain is
persistent. That is an ecological question, and experience sampling is the method built to answer it
\citep{myingermeys2018experience,shiffman2008ecological,stone1998strategies}, with an established
record in chronic pain specifically
\citep{tennen1996daily,affleck1999multilevel,peters2000electronic}.

\subsection{Two structural assumptions that are usually assumed rather than tested}
Before asking whether flow accompanies less pain, it is worth asking what a momentary flow score
actually contains. Two structural assumptions are built into the standard operationalization.

The first is configurality. The channel model does not claim that high challenge is good
or that high skill is good; it claims that their \emph{correspondence} is what matters, and that
correspondence at a high level is what produces flow. Operationally, this is usually captured by a
balance term, the negative absolute discrepancy $-|C-S|$, entered alongside an elevation term
$(C+S)/2$ that keeps a jointly low but perfectly matched moment, that is apathy, from counting as
optimal \citep{moneta1996effect}. If the configural claim is right, the balance term should carry
variance in the flow experience over and above what challenge and skill contribute on their own.
The evidence here is mixed: meta-analytically, the challenge-skill balance relates to flow only
modestly \citep{fong2015challenge}, and experimental work has found the balance effect to be
moderated rather than universal \citep{engeser2008flow}. A further problem is methodological
rather than substantive. A term built from an absolute difference imposes constraints that are
rarely examined and is only interpretable if a point on the challenge scale means the same thing
as a point on the skill scale; the recommended alternative is to keep the two components in their
own metric and to test the congruence hypothesis as a constrained case of a polynomial response
surface \citep{edwards1993polynomial,edwards1994congruence}.

The second is composite coherence. Momentary flow is typically indexed by summing or
averaging its constituent appraisals. \citet{vanderweele2022constructed} shows why a constructed
measure of this kind can mislead. The same composite coefficient is produced by three very
different worlds: one in which all constituents relate to the outcome similarly, so that the
composite is a fair summary; one in which a single constituent does all the work, so that the
composite is a diluted version of that constituent under a more ambitious name; and one in which
constituents work in opposite directions and partly cancel, so that the composite actively hides
the structure it claims to summarize. The three cannot be distinguished from the composite itself,
and the practical implication differs sharply between them, because an intervention has to know
which component to move. The remedy is to report the constituents separately as well as jointly.

Both assumptions matter for the pain question specifically. Challenge, in a population with
chronic pain, is not a neutral quantity. Demanding activity is also physically effortful, and
effort has its own cost; a literature on activity pacing, avoidance, and endurance describes
exactly this trade-off \citep{andrews2012activity,vlaeyen2000fear}. If challenge and absorption
pull in opposite directions with respect to pain, a composite that averages them will be quiet
about the most clinically consequential part of the construct.

\subsection{The present study}
This study uses a chronic-pain experience-sampling dataset in which patients answered a short
momentary diary eight times per day for two weeks. Four of the diary items map onto the flow
model: challenge and skill form the antecedent condition, and absorption and enjoyment form the
experience itself. The same prompts recorded momentary pain, the extent to which pain interfered
with what the person was doing, attention to pain, and the threat value of pain, which makes it
possible to ask the ecological counterpart of the flow-induced-analgesia question. The data are a
secondary analysis: the source study analysed these items separately and never as a flow index
\citep{viane2004acceptance}, so every flow variable here is an analyst-constructed measure and is
reported as such. Because the data are observational, all associations are read as within-person
predictive relations rather than causal effects \citep{rohrer2018thinking}. Three research questions guide the analysis.

\begin{description}[leftmargin=2.2em,itemsep=2pt]
  \item[\textbf{RQ1.}] Does the challenge-skill condition relate to the momentary flow experience
  configurationally, as the channel model requires, or additively?
  \item[\textbf{RQ2.}] Does the flow composite carry information about the momentary experience of
  pain that its four constituents, entered separately, do not?
  \item[\textbf{RQ3.}] Is the momentary flow experience associated with a better experience of pain
  in daily life, concurrently and prospectively, and is the flow channel the least painful region
  of the challenge-skill plane?
\end{description}

\section{Method}

\subsection{Participants and design}
The data are a secondary analysis of a chronic-pain experience-sampling study reported by
\citet{viane2004acceptance}, which examined acceptance of pain in relation to attention to pain
and engagement with daily activities \citep[see also][]{viane2003acceptance}; the flow question
asked here was not part of that report. Patients with chronic musculoskeletal pain, recruited
through a university pain clinic, first completed a battery of validated questionnaires and then
carried a palmtop computer that signalled them eight times per day for two weeks to complete a
short momentary diary, together with a brief morning and evening diary. Sixty-eight participants
met the study's compliance criterion of at least 25 completed prompts. Of their completed prompts,
$5{,}976$ carried all four activity appraisals and form the analytic frame for the analyses
(median $=92$ per person); the $4.6\%$ of prompts lost to a missing appraisal are characterized in
\Cref{sfig:frame} and \Cref{tab:sup-robustness}. Baseline questionnaires were available for 67
participants. Sample characteristics and the momentary appraisals are summarized in
\Cref{tab:frame}.

\begin{table}[!htbp]
\raggedright
\caption{Sample and the momentary appraisals that operationalize flow.}
\label{tab:frame}
\footnotesize
\begingroup
\setlength{\tabcolsep}{4pt}
\renewcommand{\arraystretch}{1.08}
\noindent\textit{Panel A.} Baseline characteristics of the analytic sample.\par\vspace{2pt}
\begin{tabular*}{\textwidth}{@{\extracolsep{\fill}}lccc@{}}
\toprule
Variable & Mean (SD) & Range & $n$ \\
\midrule
Age (years) & 46.4 (9.8) & 21 to 65 & 67 \\
Sex: women / men (\%) & 80.6 / 19.4 &  & 67 \\
Pain duration (months) & 82.8 (121.4) & 2 to 520 & 62 \\
MPI pain severity & 4.3 (0.9) & 2 to 6 & 67 \\
MPI interference & 4.2 (1.0) & 2 to 6 & 67 \\
Pain Disability Index & 42.7 (11.8) & 11 to 68 & 67 \\
PCS catastrophizing (total) & 22.3 (10.2) & 1 to 46 & 66 \\
PVAQ pain vigilance (total) & 43.7 (12.8) & 5 to 66 & 65 \\
HADS anxiety & 7.7 (3.9) & 1 to 20 & 66 \\
HADS depression & 8.7 (3.7) & 1 to 16 & 66 \\
NEO-FFI neuroticism & 30.8 (9.4) & 12 to 54 & 64 \\
\bottomrule
\end{tabular*}
\par\vspace{6pt}\noindent\textit{Panel B.} The four momentary appraisals and the flow experience composite.\par\vspace{2pt}
\begin{tabular*}{\textwidth}{@{\extracolsep{\fill}}llccccc@{}}
\toprule
Momentary item & Role & $n$ obs & Mean & SD & ICC & Within-person share \\
\midrule
Challenge & Condition & 6,165 & 3.28 & 1.98 & 0.242 & 0.758 \\
Skill & Condition & 6,108 & 4.79 & 1.55 & 0.226 & 0.774 \\
Absorption & Experience & 6,143 & 4.00 & 1.90 & 0.213 & 0.787 \\
Enjoyment & Experience & 6,136 & 4.90 & 1.58 & 0.207 & 0.793 \\
Flow experience & Experience & 6,035 & 4.45 & 1.51 & 0.201 & 0.799 \\
\bottomrule
\end{tabular*}
\par\vspace{3pt}\noindent\parbox{\textwidth}{\footnotesize\textit{Note.} MPI = West Haven-Yale Multidimensional Pain Inventory; PDI = Pain Disability Index; PCS = Pain Catastrophizing Scale; PVAQ = Pain Vigilance and Awareness Questionnaire; HADS = Hospital Anxiety and Depression Scale; NEO-FFI = NEO Five-Factor Inventory. Items were rated 1 (not at all) to 7 (very much). Condition items are challenge and skill, experience items absorption and enjoyment; the flow experience is their mean. ICC is the between-person variance share. The four appraisals are single momentary items; every flow variable is a secondary, analyst-constructed measure, since the source study analysed the items separately.}
\endgroup
\end{table}
}

\subsection{Momentary measures}
At each prompt, participants described the activity they were engaged in and rated it on
seven-point scales (1 = not at all, 7 = very much). Four ratings operationalize flow. Challenge was
``this activity is a personal challenge for me'' and skill was ``I am good at this''; together they
form the antecedent condition. Absorption was ``I am absorbed in this activity'' and enjoyment was
``I like doing this''; together they form the experience, absorption supplying the concentration
component and enjoyment the autotelic marker. The same prompt recorded momentary pain (``I am in
pain now''), pain interference (``my pain hinders me''), attention to pain (``I focus on my
pain''), and the threat value of pain, scored as the mean of momentary fear (``I am afraid of the
pain now'') and catastrophizing (``I feel that my pain is becoming too much for me''). A
self-reported physical activation item was recorded as well and is used as a covariate throughout,
because challenge co-varied strongly with being physically active ($r=.57$ within persons) and
activity has its own relation to pain.

These four appraisals cover three of the nine dimensions Csikszentmihalyi describes, namely the
challenge-skill configuration, concentration through absorption, and the autotelic quality through
enjoyment. Clear goals, unambiguous feedback, sense of control, loss of self-consciousness, and
time transformation were not measured. The instrument is therefore a core flow proxy rather than a
flow scale \citep{jackson2002assessing}, and it is described as such throughout.

\subsection{Baseline measures}
Baseline questionnaires provided three person-level profiles used in the person-level analyses.
The threat-value profile combined the Pain Catastrophizing Scale total \citep{sullivan1995pain} and
the Pain Vigilance and Awareness Questionnaire total \citep{roelofs2003pvaq}. The biomedical
profile combined the West Haven-Yale Multidimensional Pain Inventory pain-severity subscale
\citep{kerns1985mpi} and pain duration. The personality profile was NEO-FFI neuroticism
\citep{costa1992neo}. Each profile was standardized.

\subsection{Operationalizing flow}
Following the standard experience-sampling operationalization
\citep{massimini1988systematic,moneta1996effect}, the four items were represented as two
structurally distinct parts rather than as one sum, because summing them conflates the antecedent
condition with the outcome experience. The condition was expressed as balance and elevation, and
the experience as the mean of absorption and enjoyment:
%
\begin{equation}
\text{balance} = -\lvert C - S \rvert, \qquad
\text{elevation} = \tfrac{1}{2}(C + S), \qquad
\text{flow experience} = \tfrac{1}{2}(A + E),
\label{eq:flow}
\end{equation}
%
with $C$ challenge, $S$ skill, $A$ absorption, and $E$ enjoyment. The elevation term is not
optional: without it, a moment low on both challenge and skill would score as perfectly balanced
and therefore as optimal, which is apathy rather than flow. For the descriptive channel analysis a
moment was assigned to flow when both condition items were high, to anxiety when only challenge
was high, to relaxation when only skill was high, and to apathy when neither was
\citep{carli1988quality}.

Two measurement decisions shape everything that follows and are reported explicitly. The first is
the metric. Person-mean centering was primary \citep{curran2011disaggregation}, because it
separates a person's momentary fluctuation from their level. Within-person standardization was
run as a sensitivity analysis but cannot carry prevalence claims: it sets every participant's mean
to zero and therefore assigns above-average moments to everyone, including a participant who never
reaches flow in any absolute sense. All prevalence statements are therefore made on the raw
absolute scale, with both a conservative criterion (both condition items at least 5) and a liberal
one (both at least 4); the grand-mean $z$ metric was run as a further sensitivity analysis
(\Cref{sfig:standardization}, \Cref{tab:sup-prevalence}, \Cref{tab:sup-standardization}). The
second is provenance: these are single momentary items, not validated subscales, and the source
study analysed them separately, so every flow variable is a secondary, analyst-constructed measure.

\subsection{Statistical modelling}
Analyses were implemented in a reproducible pipeline in which statistical models were written in
\textsf{R} \citep{rcoreteam2024} and orchestrated from \textsf{Python}, which handled
preprocessing, tables, and visualization. Multilevel models were fitted with \textsf{lme4}
\citep{bates2015lme4}. Every reported table and figure was regenerated from the processed data.

\paragraph{RQ1: configural or additive.} The condition-to-experience model regressed the flow
experience on balance and elevation, decomposed into within-person and between-person components,
with random slopes for both within-person terms. Because the balance parameterization is only one
of several ways to express a configural claim, and because a term built from an absolute
difference is interpretable only under a commensurability assumption that two single items cannot
guarantee \citep{edwards1993polynomial,edwards1994congruence}, the configural hypothesis was
tested in four further forms, each of which keeps challenge and skill in their own metric and each
of which is a nested extension of the additive model: a balance term added to the additive model,
a one-sided kink at $C=S$ that allows overload to differ from boredom, a multiplicative
$C \times S$ interaction, a full second-order response surface, and the four-channel quadrant
classification. These nested comparisons were estimated by maximum likelihood with a random
intercept, so that only the fixed-effect structure differs between the compared models, and were
evaluated by likelihood-ratio test, AIC, and BIC. The two parameterizations of the condition,
balance-and-elevation versus additive challenge-and-skill, have the same number of parameters and
are not nested, so they were compared by information criterion.

\paragraph{RQ2: the composite against its constituents.} Absorption, enjoyment, challenge, and
skill were entered as the only appraisal in the model (``alone'') and then all four together
(``jointly''), predicting each of the four momentary pain measures, with physical activation as a
covariate. The comparison of the two sets is the constituent decomposition that
\citet{vanderweele2022constructed} recommends, and a change of sign between them is the signature
of statistical suppression \citep{mackinnon2000equivalence}. Because the flow experience is the
mean of absorption and enjoyment, the composite model is exactly the constituent model under the
constraint that those two carry equal weight and that challenge and skill carry none, so the two
are nested and can be compared by likelihood-ratio test. As an independent, multivariate check, a
four-node multilevel vector autoregressive network \citep{epskamp2018gaussian} of momentary pain,
its threat value, attention to pain, and absorption was estimated and then re-estimated with the
flow composite substituted for the single absorption node, and the two networks were compared edge
by edge.

\paragraph{RQ3: flow and the momentary experience of pain.} Concurrent multilevel models regressed
each pain measure on the flow experience, with the condition terms in the model, a random slope for
the flow experience, and both without and with physical activation as a covariate. Temporal models
regressed each pain measure at $t$ on the flow experience at $t-1$ while controlling the outcome's
own autoregression, and the reverse direction as well, so that the temporal picture is reported in
both directions rather than in one. Lags were built between adjacent prompts within the same study
day, so no lag crosses the overnight gap, and were estimated on person-standardized scores. The
four channels were compared with apathy as the reference, the conventional contrast, and again with
flow as the reference, which turns the ordering claim into a direct test. Person-specific
associations were estimated by unregularized ordinary least squares for each participant with at
least 50 usable moments.

\paragraph{Robustness.} Every focal model was re-estimated on all four standardization metrics,
with day and prompt position as level-1 covariates, and at compliance floors of 20, 30, 50, and 70
completed moments per person. Moments lost to listwise deletion were compared with retained moments
on every momentary measure. No correction for multiplicity was applied; the nested tests are
pre-specified contrasts of a single construct rather than a screen.

\subsection{Hypotheses}
Six hypotheses follow from the flow literature and from the flow-induced-analgesia proposal. For
RQ1, H1 predicted configurality: balance should relate positively to the flow experience over and
above elevation, and more generally some configural term should improve on the additive model. For
RQ2, H2 predicted composite coherence: the four constituents should relate to pain in the same
direction and to a comparable degree, so that entering them separately adds nothing. For RQ3, H3
predicted a concurrent association, with the flow experience accompanying lower pain and, more
strongly, less pain interference, since an attentional account targets the interruptive rather than
the sensory aspect of pain; H4 predicted the temporal counterpart, that a higher than usual flow
experience would be followed by lower pain at the next prompt; and H5 predicted that the flow
channel would be the least painful of the four channels. H6 predicted that these associations would
vary substantially across patients.

\section{Results}

\subsection{The instrument, and how often it finds flow}
The four appraisals behaved as momentary states rather than as traits: between $76$ and $79\%$ of
each item's variance was within person (ICCs $.21$ to $.24$; \Cref{fig:instrument}, Panel C,
\Cref{tab:frame}). Their levels were not symmetric. Skill (mean $=4.79$) and enjoyment
(mean $=4.90$) sat well above the scale midpoint while challenge sat well below it
(mean $=3.28$), so most everyday activity in this sample was manageable and liked but not
demanding (\Cref{fig:instrument}, Panel A). Within persons, challenge correlated only $r=.17$ with
skill, whereas absorption correlated $r=.53$ with challenge and $r=.42$ with skill
(\Cref{fig:instrument}, Panel B, \Cref{tab:sup-itemcorr}).

How often the diary calls a moment a flow moment depended almost entirely on the classification
rule (\Cref{fig:instrument}, Panel D, \Cref{tab:sup-prevalence}). Under the conservative absolute
criterion, $19.3\%$ of moments qualified and four of 68 participants never reached the criterion at
all; under the liberal absolute criterion $36.8\%$ qualified and no participant scored zero. The
relative rules gave intermediate figures ($24.6$ to $27.8\%$) but, as expected, placed no
participant at zero, which is why they cannot carry a prevalence claim. Person-level proneness
ranged from never to two thirds of a participant's moments (\Cref{fig:instrument}, Panel E). One
descriptive result is worth stating before any modelling: under the conservative rule the modal
channel was not flow ($22.0\%$ of moments met the condition) but relaxation, that is high skill with
low challenge ($37.0\%$), with apathy next ($32.3\%$) and anxiety rare ($8.6\%$;
\Cref{fig:instrument}, Panel F).

\begin{figure}[htbp]
\raggedright
\caption{The four momentary appraisals and how often the diary calls a moment a flow moment.}
\label{fig:instrument}
\noindent\includegraphics[width=0.98\textwidth]{\asset{figures/main/MAIN_01_flow_instrument.png}}
\figurenote{Panel A shows the response distribution of the four appraisals, labelled by their role
in the flow model. Panel B shows their correlation structure, within persons below the diagonal and
between persons above it. Panel C shows the share of each item's variance that lies within persons.
Panel D shows the share of moments classified as flow under each rule, before and after the
experience gate, annotated with the number of participants who never reach the criterion. Panel E
shows the per-person share of flow moments under the liberal and the strict absolute criterion.
Panel F shows the share of moments in each channel under the strict rule, annotated with the mean
flow experience of that channel.}
\end{figure}

\subsection{RQ1: the condition acts additively, not configurationally}
In the standard parameterization, the challenge-skill condition predicted the flow experience
strongly, and did so entirely through elevation (within persons $b=.744$, $SE=.024$, $p<.001$;
between persons $b=.796$, $SE=.085$, $p<.001$) rather than through balance (within persons
$b=.006$, $SE=.020$, $95\%$ CI $[-.033, .045]$, $p=.78$; \Cref{fig:configural}, Panel A,
\Cref{tab:condition}). With just under $6{,}000$ moments this is an informative null rather than an
absence of power: the interval rules out any balance effect larger than $.05$, against an elevation
effect of $.74$ on the same scale \citep{lakens2018equivalence}. Between persons the balance effect
was not merely absent but reversed, so that participants whose activities were on average better
matched reported on average \emph{less} flow experience ($b=-.111$, $SE=.049$, $p=.028$). The
reversal was not an artefact of centering: on the raw response scale it was larger and more
precisely estimated ($b=-.090$, $p<.001$), and across all four standardization metrics the
elevation effect stayed between $.74$ and $.89$ while the balance effect was never positive
(\Cref{fig:configural}, Panel B, \Cref{tab:sup-standardization}). Person by person, the balance
slope was centred on zero (median $=+.009$, range $-.27$ to $+.31$) whereas the elevation slope was
positive in all 68 participants (median $=+.758$, range $+.23$ to $+.98$; \Cref{fig:configural},
Panel F).

Because the balance term expresses only one of several possible configural claims, and because it
is interpretable only if a point on the challenge scale means the same as a point on the skill
scale \citep{edwards1993polynomial}, the configural hypothesis was tested in four further forms,
each keeping challenge and skill in their own metric and each nested in a comparison with the
additive model (\Cref{fig:configural}, Panel C, \Cref{tab:condition}). None improved on additivity.
Adding balance to a model that already contained challenge and skill gained nothing
($\chi^2(1)=0.56$, $p=.45$, $\Delta$AIC $=-1.4$). Allowing a one-sided kink at $C=S$, so that
overload could differ from boredom, gained nothing ($\chi^2(1)=0.56$, $p=.45$). A multiplicative
$C \times S$ interaction gained nothing ($\chi^2(1)=0.00$, $p=.97$), as it also did not in the
random-slope model reported in \Cref{tab:sup-condmodels} ($p=.15$). A full second-order response
surface gained nothing ($\chi^2(3)=3.93$, $p=.27$). The four-channel quadrant classification was
the only configural form to reach conventional significance ($\chi^2(3)=8.16$, $p=.043$,
$\Delta$AIC $=+2.2$), but BIC preferred the additive model over it by $17.9$, so the quadrant
scheme buys a marginal improvement at a cost that the more conservative criterion rejects. The
additive model also decisively outperformed the balance-and-elevation form at the same number of
parameters ($\Delta$AIC $=265.7$ with a random intercept, and $441.1$ in the random-slope models
of \Cref{tab:condition}), which is what should happen when a parameterization discards information:
the balance form constrains challenge and skill to equal weight, whereas the free additive estimates
put skill at roughly twice challenge ($b=.507$ versus $b=.285$; \Cref{fig:configural}, Panel E).

Two comparisons in which balance does appear to matter are reported for completeness and are
diagnostics of the parameterization rather than evidence for configurality
(\Cref{sfig:baldiag}, \Cref{tab:sup-baldiag}). Balance interacts with elevation inside the
balance-and-elevation form ($\chi^2(1)=17.6$, $p<.001$), with a balance effect that is negative at
low elevation ($b=-.031$, $p=.06$) and positive at high elevation ($b=+.070$, $p<.001$); this is
what happens when a parameterization forces two components with unequal weights into a single
symmetric difference, and it disappears once challenge and skill are free. Balance also gains
significance when added to the full quadratic surface ($\chi^2(1)=18.87$, $p<.001$), but balance
correlates $.93$ with the squared discrepancy that the surface already contains, so neither
coefficient is interpretable alone, and the joint test of the surface against additivity is null.
The most direct evidence is descriptive and needs no parameterization at all: within a band of
elevation, mean flow experience is flat or slightly declining across the whole range of mismatch
(\Cref{fig:configural}, Panel D), and the residuals of the additive model show no ridge along the
$C=S$ diagonal (\Cref{sfig:baldiag}, Panel A). H1 was therefore not supported.

\begin{figure}[htbp]
\raggedright
\caption{Configural and additive forms of the challenge-skill condition.}
\label{fig:configural}
\noindent\includegraphics[width=0.98\textwidth]{\asset{figures/main/MAIN_02_configural_versus_additive.png}}
\figurenote{Panel A shows the balance and elevation effects on the flow experience, within and
between persons. Panel B shows the same two terms estimated on each of the four standardization
metrics. Panel C shows the AIC gained over the additive challenge-plus-skill model by each
configural extension, annotated with its likelihood-ratio test; green marks $p<.05$. Panel D shows
the observed mean flow experience against the size of the challenge-skill mismatch within three
bands of elevation. Panel E shows the additive main effects and the multiplicative interaction.
Panel F shows the person-specific balance and elevation slopes. Error bars and intervals are
$95\%$ confidence intervals.}
\end{figure}

\begin{table}[!htbp]
\raggedright
\caption{The challenge-skill condition predicting the flow experience, and nested tests of configurality.}
\label{tab:condition}
\footnotesize
\begingroup
\setlength{\tabcolsep}{4pt}
\renewcommand{\arraystretch}{1.08}
\noindent\textit{Panel A.} Balance and elevation predicting the flow experience (multilevel model with random slopes).\par\vspace{2pt}
\begin{tabular*}{\textwidth}{@{\extracolsep{\fill}}lccccc@{}}
\toprule
Term & $b$ & SE & 95\% CI & $p$ & Random SD \\
\midrule
Balance, within person & 0.006 & 0.020 & [-0.033, 0.045] & .782 & 0.131 \\
Elevation, within person & 0.744 & 0.024 & [0.697, 0.791] & <.001 & 0.164 \\
Balance, between persons & -0.111 & 0.049 & [-0.207, -0.015] & .028 &  \\
Elevation, between persons & 0.796 & 0.085 & [0.629, 0.963] & <.001 &  \\
\bottomrule
\end{tabular*}
\par\vspace{6pt}\noindent\textit{Panel B.} Nested likelihood-ratio tests of every configural functional form, each added to the additive challenge-plus-skill model.\par\vspace{2pt}
\begin{tabular*}{\textwidth}{@{\extracolsep{\fill}}lccccc@{}}
\toprule
Term added to challenge + skill & $\Delta$ df & $\chi^2$ & $p$ & $\Delta$ AIC & $\Delta$ BIC \\
\midrule
Balance, $-|C-S|$ & 1 & 0.56 & .452 & -1.4 & -8.1 \\
One-sided kink at $C=S$ & 1 & 0.56 & .452 & -1.4 & -8.1 \\
Interaction, $C \times S$ & 1 & 0.00 & .970 & -2.0 & -8.7 \\
Quadratic response surface & 3 & 3.93 & .270 & -2.1 & -22.2 \\
Four-channel quadrant & 3 & 8.16 & .043 & 2.2 & -17.9 \\
\bottomrule
\end{tabular*}
\par\vspace{6pt}\noindent\textit{Panel C.} The two parameterizations of the condition at equal degrees of freedom.\par\vspace{2pt}
\begin{tabular*}{\textwidth}{@{\extracolsep{\fill}}lcccc@{}}
\toprule
Parameterization & df & AIC & BIC & $\Delta$ AIC \\
\midrule
balance + elevation & 5 & 17545.4 & 17578.8 & 265.7 \\
challenge + skill (additive) & 5 & 17279.7 & 17313.2 & 0.0 \\
\bottomrule
\end{tabular*}
\par\vspace{3pt}\noindent\parbox{\textwidth}{\footnotesize\textit{Note.} Balance is $-|C-S|$ and elevation $(C+S)/2$, both person-mean centred; $C$ is challenge and $S$ skill. Panel A is estimated with REML and random slopes for both within-person terms. Panels B and C are estimated with maximum likelihood and a random intercept, so that only the fixed-effect structure differs across the compared models. Positive $\Delta$ AIC and $\Delta$ BIC favour the more complex model. The four appraisals are single momentary items; every flow variable is a secondary, analyst-constructed measure, since the source study analysed the items separately.}
\endgroup
\end{table}
}

\subsection{RQ2: the composite hides a sign reversal among its constituents}
Entered as the only appraisal in the model, three of the four constituents tracked lower momentary
pain: skill ($b=-.088$), enjoyment ($b=-.077$), and absorption ($b=-.069$, all $p<.001$), while
challenge was flat ($b=-.004$, $p=.82$). The composite behaved as their average
($b=-.096$, $p<.001$). Entered jointly, the picture changed in three ways
(\Cref{fig:decomposition}, Panels A and B, \Cref{tab:decomposition}). Absorption ($b=-.055$,
$p<.001$) and skill ($b=-.045$, $p<.001$) kept their associations. Enjoyment fell to
non-significance ($b=-.023$, $p=.089$), so its apparent relation to pain was variance it shared
with the other two rather than a contribution of its own. And challenge reversed sign and became
reliably positive ($b=+.032$, $p=.004$): once skill and absorption are held constant, a more
demanding moment is a more painful one. The reversal is the classic signature of suppression
\citep{mackinnon2000equivalence}, and its source is visible in the correlation structure, since
challenge co-varies $r=.53$ with absorption and $r=.17$ with skill, both of which track lower pain
(\Cref{fig:decomposition}, Panel E). It was not confined to pain intensity: challenge was reliably
positive for pain interference ($b=+.053$), attention to pain ($b=+.031$), and the threat value of
pain ($b=+.035$, all $p<.01$), so on every outcome one constituent of the flow construct works
against the other three (\Cref{fig:decomposition}, Panel B).

Because the flow experience is the mean of absorption and enjoyment, the composite model is the
constituent model under a linear constraint, and the two can be compared directly. For pain
intensity, splitting the composite into its two experience components changed nothing
($\chi^2(1)=0.02$, $p=.89$), so the composite is a fair summary of \emph{those two}. Adding the two
condition items, however, improved fit substantially for every outcome (pain $\chi^2(3)=21.0$,
interference $\chi^2(3)=36.1$, attention $\chi^2(3)=26.5$, threat $\chi^2(3)=34.0$, all $p<.001$;
\Cref{fig:decomposition}, Panel D). The flow architecture as a whole therefore loses information
relative to its parts, and it loses it in the specific way \citet{vanderweele2022constructed}
warns about, by averaging over constituents that pull in opposite directions.

An independent, multivariate check pointed the same way. Substituting the flow composite for the
single absorption node in a four-node temporal network of momentary pain, its threat value,
attention to pain, and absorption reproduced that network almost exactly: across all 22 estimated
parameters the two networks correlated $r=.997$, the
largest edge change was $.027$, and no edge changed sign or significance
(\Cref{fig:decomposition}, Panel F, \Cref{tab:sup-network}). In a dynamic system of four nodes,
then, the composite behaves indistinguishably from its dominant constituent. H2 was not supported.

\begin{figure}[htbp]
\raggedright
\caption{The flow composite against its four constituents.}
\label{fig:decomposition}
\noindent\includegraphics[width=0.98\textwidth]{\asset{figures/main/MAIN_03_constituent_decomposition.png}}
\figurenote{Panel A shows each constituent's association with momentary pain when it is the only
appraisal in the model (hollow) and when all four are entered together (filled), with the composite
marked. Panel B shows the joint coefficients for all four pain measures. Panel C shows how far each
coefficient moves between the two models. Panel D shows the AIC gained over the composite by
splitting it in two and by using all four constituents. Panel E shows the within-person
correlations among the constituents that produce the suppression. Panel F plots every parameter of
the benchmark temporal and contemporaneous networks against the same parameter after the composite
is substituted for the absorption node. Error bars are $95\%$ confidence intervals.}
\end{figure}

\begin{table}[!htbp]
\raggedright
\caption{Constituent decomposition of the flow composite.}
\label{tab:decomposition}
\footnotesize
\begingroup
\setlength{\tabcolsep}{4pt}
\renewcommand{\arraystretch}{1.08}
\noindent\textit{Panel A.} Each appraisal entered alone, predicting the four momentary pain measures.\par\vspace{2pt}
\begin{tabular*}{\textwidth}{@{\extracolsep{\fill}}lcccc@{}}
\toprule
Constituent & Pain intensity & Pain interference & Attention to pain & Threat value \\
\midrule
Absorption & -0.069*** (0.014) & -0.089*** (0.017) & -0.048*** (0.013) & -0.039*** (0.009) \\
Enjoyment & -0.077*** (0.016) & -0.118*** (0.021) & -0.074*** (0.015) & -0.066*** (0.012) \\
Challenge & -0.004 (0.017) & +0.004 (0.022) & +0.006 (0.015) & +0.017 (0.012) \\
Skill & -0.088*** (0.016) & -0.125*** (0.021) & -0.086*** (0.017) & -0.059*** (0.013) \\
Flow experience (composite) & -0.096*** (0.018) & -0.133*** (0.022) & -0.077*** (0.016) & -0.068*** (0.012) \\
\bottomrule
\end{tabular*}
\par\vspace{6pt}\noindent\textit{Panel B.} The same four appraisals entered jointly.\par\vspace{2pt}
\begin{tabular*}{\textwidth}{@{\extracolsep{\fill}}lcccc@{}}
\toprule
Constituent & Pain intensity & Pain interference & Attention to pain & Threat value \\
\midrule
Absorption & -0.055*** (0.011) & -0.062*** (0.014) & -0.029* (0.011) & -0.023** (0.009) \\
Enjoyment & -0.023 (0.013) & -0.057*** (0.016) & -0.035** (0.013) & -0.044*** (0.010) \\
Challenge & +0.032** (0.011) & +0.053*** (0.013) & +0.031** (0.011) & +0.035*** (0.008) \\
Skill & -0.045*** (0.013) & -0.063*** (0.016) & -0.049*** (0.013) & -0.026* (0.010) \\
\bottomrule
\end{tabular*}
\par\vspace{6pt}\noindent\textit{Panel C.} Nested likelihood-ratio tests of the composite against its constituents.\par\vspace{2pt}
\begin{tabular*}{\textwidth}{@{\extracolsep{\fill}}llcccc@{}}
\toprule
Outcome & Comparison & $\Delta$ df & $\chi^2$ & $p$ & $\Delta$ AIC \\
\midrule
Pain intensity & Composite vs. absorption + enjoyment & 1 & 0.02 & .891 & -2.0 \\
Pain intensity & Absorption + enjoyment vs. all four & 2 & 20.98 & <.001 & 17.0 \\
Pain intensity & Composite vs. all four constituents & 3 & 21.00 & <.001 & 15.0 \\
Pain interference & Composite vs. absorption + enjoyment & 1 & 2.64 & .104 & 0.6 \\
Pain interference & Absorption + enjoyment vs. all four & 2 & 33.49 & <.001 & 29.5 \\
Pain interference & Composite vs. all four constituents & 3 & 36.13 & <.001 & 30.1 \\
Attention to pain & Composite vs. absorption + enjoyment & 1 & 3.49 & .062 & 1.5 \\
Attention to pain & Absorption + enjoyment vs. all four & 2 & 22.96 & <.001 & 19.0 \\
Attention to pain & Composite vs. all four constituents & 3 & 26.45 & <.001 & 20.5 \\
Threat value & Composite vs. absorption + enjoyment & 1 & 8.74 & .003 & 6.7 \\
Threat value & Absorption + enjoyment vs. all four & 2 & 25.30 & <.001 & 21.3 \\
Threat value & Composite vs. all four constituents & 3 & 34.04 & <.001 & 28.0 \\
\bottomrule
\end{tabular*}
\par\vspace{3pt}\noindent\parbox{\textwidth}{\footnotesize\textit{Note.} Within-person (person-mean centred) predictors and outcomes. In Panel A ``alone'' means the constituent is the only appraisal in the model and ``joint'' means all four are entered together; the bottom row gives the composite for reference. $^{*}p<.05$, $^{**}p<.01$, $^{***}p<.001$. Panel B compares nested fixed-effect structures under maximum likelihood with a random intercept and momentary physical activation as a covariate; the composite is the constituent model under the constraint that absorption and enjoyment carry equal weight and that challenge and skill carry none. The four appraisals are single momentary items; every flow variable is a secondary, analyst-constructed measure, since the source study analysed the items separately.}
\endgroup
\end{table}
}

\subsection{RQ3: a concurrent association, no temporal one, and a channel ordering that
inverts the prediction}
At the same prompt, a higher than usual flow experience accompanied a better experience of pain on
all four measures: less pain interference ($b=-.139$, $SE=.025$), lower pain intensity
($b=-.094$, $SE=.020$), a lower threat value ($b=-.079$, $SE=.014$), and less attention to pain
($b=-.074$, $SE=.019$; all $p<.001$; \Cref{fig:flowpain}, Panel A, \Cref{tab:flowpain}). All four
survived adjustment for physical activation, which matters because challenge co-varies with being
active and activity independently raises pain in this sample (\Cref{sfig:robust}, Panels A and B).
The ordering of the effects is itself informative and is what an attentional account predicts:
interference, the interruptive aspect of pain, moved most, and moved half again as much as
intensity, the sensory aspect. H3 was supported.

The temporal test gave the opposite answer. No lag-1 effect approached significance in either
direction: neither the flow experience at $t-1$ predicting any pain measure at $t$
($p$ between $.44$ and $.92$), nor any pain measure at $t-1$ predicting the flow experience at $t$
($p$ between $.28$ and $.94$), with every point estimate below $|b|=.02$ (\Cref{fig:flowpain},
Panel B). This is not a failure of the design to see time. In the same models the autoregressive
effects were large and highly significant (pain $b=.395$, interference $b=.352$, threat $b=.321$,
attention $b=.266$, all $p<.001$), so the models detect temporal structure where it exists and
detect none between flow and pain (\Cref{sfig:temporal}, Panel D). The same held inside the flow
construct itself: the challenge-skill condition at $t-1$ did not predict the flow experience at
$t$ (elevation $b=.052$, $p=.062$; balance $b=-.007$, $p=.70$), which is $6\%$ of the same-prompt
elevation effect estimated on the identical rows ($b=.856$; \Cref{sfig:temporal}, Panels A and B,
\Cref{tab:sup-condlag}). The flow-pain relation in these data is co-occurrence, not prediction, and
without temporal precedence the direction cannot be identified: absorption may dampen pain, or low
pain may make absorption possible. H4 was not supported.

The channel comparison inverted the predicted ordering. Against apathy, the conventional reference,
flow looked good on every measure (pain $b=-.179$, interference $b=-.173$, attention $b=-.169$,
threat $b=-.070$), but relaxation looked better still (pain $b=-.210$, interference $b=-.272$,
attention $b=-.280$, threat $b=-.167$; \Cref{fig:plane}, Panel D, \Cref{tab:channels}). Tested
head to head with flow as the reference, relaxation was reliably better on pain interference
($b=-.098$, $p=.041$), attention to pain ($b=-.111$, $p=.005$), and the threat value of pain
($b=-.097$, $p=.001$), and indistinguishable from flow on pain intensity ($b=-.032$, $p=.42$);
anxiety and apathy were reliably worse than flow on all four (\Cref{fig:plane}, Panel E). Plotted
over the full challenge-skill plane, the least painful region was the high-skill, low-challenge
corner rather than the high-high corner (\Cref{fig:plane}, Panels B and C). This follows directly
from the decomposition: within persons, skill tracked lower pain while challenge did not, and once
entered jointly challenge tracked higher pain, so the most favourable corner of the plane is
low-challenge and high-skill by construction (\Cref{fig:plane}, Panel F). H5 was therefore
rejected in favour of the relaxation channel on three of four outcomes and was not supported on the
fourth.

\begin{figure}[htbp]
\raggedright
\caption{Concurrent, lagged, and person-specific associations between the flow experience and pain.}
\label{fig:flowpain}
\noindent\includegraphics[width=0.98\textwidth]{\asset{figures/main/MAIN_04_flow_and_pain.png}}
\figurenote{Panel A shows the concurrent association of the flow experience with the four momentary
pain measures, unadjusted and adjusted for physical activation. Panel B shows every lag-1 effect in
both directions, with the four autoregressive effects from the same models below the dotted line as
the reference for what a detectable temporal effect looks like here. Panels C and D show the
person-specific concurrent and lag-1 slopes with their $95\%$ confidence intervals, ordered by size;
intervals excluding zero are highlighted. Panel E overlays the two person-level distributions.
Panel F shows the concurrent estimate at four compliance floors.}
\end{figure}

\begin{table}[!htbp]
\raggedright
\caption{The flow experience and the momentary experience of pain.}
\label{tab:flowpain}
\footnotesize
\begingroup
\setlength{\tabcolsep}{4pt}
\renewcommand{\arraystretch}{1.08}
\noindent\textit{Panel A.} The flow experience and the momentary experience of pain at the same beep.\par\vspace{2pt}
\begin{tabular*}{\textwidth}{@{\extracolsep{\fill}}lccccccc@{}}
\toprule
Outcome & $b$ unadjusted & SE & $b$ adjusted & SE  & 95\% CI & $p$ & Random SD \\
\midrule
Pain intensity & -0.096 & 0.020 & -0.094 & 0.020 & [-0.133, -0.055] & <.001 & 0.114 \\
Pain interference & -0.143 & 0.025 & -0.139 & 0.025 & [-0.188, -0.090] & <.001 & 0.146 \\
Attention to pain & -0.077 & 0.019 & -0.074 & 0.019 & [-0.111, -0.037] & <.001 & 0.097 \\
Threat value & -0.080 & 0.014 & -0.079 & 0.014 & [-0.106, -0.052] & <.001 & 0.071 \\
\bottomrule
\end{tabular*}
\par\vspace{6pt}\noindent\textit{Panel B.} Lag-1 effects in both directions, with the outcome's own autoregression in the same model.\par\vspace{2pt}
\begin{tabular*}{\textwidth}{@{\extracolsep{\fill}}lcccc@{}}
\toprule
Path & $b$ & SE & $p$ & Autoregression of the outcome \\
\midrule
Flow $\rightarrow$ Pain intensity & -0.002 & 0.014 & .888 & 0.395 \\
Flow $\rightarrow$ Pain interference & 0.001 & 0.014 & .921 & 0.352 \\
Flow $\rightarrow$ Attention to pain & -0.012 & 0.016 & .451 & 0.266 \\
Flow $\rightarrow$ Threat value & -0.011 & 0.014 & .439 & 0.321 \\
Pain intensity $\rightarrow$ flow & -0.019 & 0.018 & .280 & 0.128 \\
Attention to pain $\rightarrow$ flow & -0.001 & 0.016 & .937 & 0.131 \\
Threat value $\rightarrow$ flow & -0.017 & 0.018 & .342 & 0.129 \\
\bottomrule
\end{tabular*}
\par\vspace{6pt}\noindent\textit{Panel C.} Distribution of the person-specific flow-pain slope; ``rel.'' marks slopes whose $95\%$ interval excludes zero.\par\vspace{2pt}
\begin{tabular*}{\textwidth}{@{\extracolsep{\fill}}lccccccc@{}}
\toprule
Slope & $n$ & Median & IQR & Range & \% neg. & Rel. neg. & Rel. pos. \\
\midrule
Same beep & 63 & -0.137 & [-0.20, 0.01] & [-0.47, 0.45] & 73.0 & 14 & 4 \\
Lag 1 & 63 & 0.000 & [-0.08, 0.05] & [-0.27, 0.23] & 49.2 & 2 & 0 \\
\bottomrule
\end{tabular*}
\par\vspace{3pt}\noindent\parbox{\textwidth}{\footnotesize\textit{Note.} Panel A models are multilevel with a random slope for the flow experience and adjust for the challenge-skill condition; the adjusted column adds momentary physical activation. Panel B is estimated on person-standardized scores with lags built within day, so no lag crosses the overnight gap. Panel C reports unregularized per-person ordinary least squares estimates for participants with at least 50 usable moments; ``reliably'' means the 95\% confidence interval excludes zero. The four appraisals are single momentary items; every flow variable is a secondary, analyst-constructed measure, since the source study analysed the items separately.}
\endgroup
\end{table}
}

\subsection{Person-specific heterogeneity}
The average concurrent association summarized wide variation, supporting H6. The person-specific
flow-pain slope was negative in $73\%$ of the 63 participants with enough moments to estimate it,
with a median of $-.137$ and a range from $-.465$ to $+.445$; it was reliably negative in 14
participants and reliably positive in 4 (\Cref{fig:flowpain}, Panels C and E, \Cref{tab:flowpain}).
Those four are not noise around a common effect but participants whose confidence intervals exclude
zero on the other side, and for whom a flow-based intervention would plausibly be
counterproductive. The lag-1 slopes, by contrast, were centred on zero (median $=.000$, $49\%$
negative, reliably negative in 2 and reliably positive in none), reproducing the pooled null at the
individual level (\Cref{fig:flowpain}, Panel D). Between-person spread was of the same order as the
average effect itself: the random-slope standard deviation was $.114$ for pain intensity against a
fixed effect of $-.094$, and $.146$ for interference against $-.139$ (\Cref{sfig:robust}, Panel E).

Person-level flow variables related to the baseline profiles in a way that did not extend to the
coupling (\Cref{sfig:personlevel}, \Cref{tab:sup-personlevel}). A stronger baseline threat-value
profile went with a lower average flow experience ($r=-.27$, $p=.028$; partial $r=-.27$,
$p=.025$ controlling mean pain), as did higher neuroticism ($r=-.26$, $p=.035$; partial $r=-.29$,
$p=.022$), whereas the biomedical profile was unrelated ($r=.03$, $p=.80$). No profile predicted the
person-specific flow-pain slope (all $p \geq .18$). Baseline characteristics thus indexed how much
flow a patient typically experienced but not how that experience related to their pain, so the
heterogeneity in the coupling is not reducible to the trait differences measured here.

\subsection{Robustness}
The conclusions did not depend on the analytic choices that could have driven them
(\Cref{sfig:robust}, \Cref{tab:sup-robustness}). The concurrent flow-pain estimate moved from
$-.094$ to $-.088$ across compliance floors of 20, 30, 50, and 70 completed moments per person and
remained significant throughout. Adding day and prompt position as level-1 covariates left every
focal estimate essentially unchanged (elevation $.744$ to $.746$; flow to pain $-.094$ to $-.101$),
although prompt position was itself a significant predictor. The condition results held on all four
standardization metrics. Listwise deletion removed $4.6\%$ of prompts; the removed moments carried
slightly more negative affect ($d=0.19$) and threat ($d=0.15$) and slightly less physical
activation ($d=-0.24$), but did not differ in pain intensity ($d=-0.12$, $p=.05$),
interference, or attention (\Cref{sfig:frame}, Panel D).

\begin{figure}[htbp]
\raggedright
\caption{The challenge-skill plane and where the least painful moments sit.}
\label{fig:plane}
\noindent\includegraphics[width=0.98\textwidth]{\asset{figures/main/MAIN_05_challenge_skill_plane.png}}
\figurenote{Panels A to C show the mean flow experience, mean pain intensity, and mean pain
interference across the full raw challenge-skill plane, with the four channels marked and cells
holding fewer than 15 moments left blank. Panel D shows each channel against apathy, the
conventional reference. Panel E shows each channel against flow, the direct head-to-head test;
grey intervals are not significant. Panel F shows the person-centred mean pain at each level of
challenge and of skill. Error bars and intervals are $95\%$ confidence intervals.}
\end{figure}

\begin{table}[!htbp]
\raggedright
\caption{The four challenge-skill channels and the momentary experience of pain.}
\label{tab:channels}
\footnotesize
\begingroup
\setlength{\tabcolsep}{4pt}
\renewcommand{\arraystretch}{1.08}
\noindent\textit{Panel A.} Each channel against apathy, the conventional reference.\par\vspace{2pt}
\begin{tabular*}{\textwidth}{@{\extracolsep{\fill}}lcccc@{}}
\toprule
Channel & Pain intensity & Pain interference & Attention to pain & Threat value \\
\midrule
Flow & -0.179 (<.001) & -0.173 (<.001) & -0.169 (<.001) & -0.070 (.018) \\
Relaxation & -0.210 (<.001) & -0.272 (<.001) & -0.280 (<.001) & -0.167 (<.001) \\
Anxiety & +0.028 (.522) & +0.131 (.015) & -0.006 (.888) & +0.043 (.201) \\
\bottomrule
\end{tabular*}
\par\vspace{6pt}\noindent\textit{Panel B.} Each channel against flow, the direct head-to-head comparison.\par\vspace{2pt}
\begin{tabular*}{\textwidth}{@{\extracolsep{\fill}}lcccc@{}}
\toprule
Channel & Pain intensity & Pain interference & Attention to pain & Threat value \\
\midrule
Relaxation & -0.032 (.425) & -0.098 (.041) & -0.111 (.005) & -0.097 (.001) \\
Anxiety & +0.207 (<.001) & +0.305 (<.001) & +0.162 (<.001) & +0.113 (<.001) \\
Apathy & +0.179 (<.001) & +0.173 (<.001) & +0.169 (<.001) & +0.070 (.018) \\
\bottomrule
\end{tabular*}
\par\vspace{3pt}\noindent\parbox{\textwidth}{\footnotesize\textit{Note.} Channels are formed from person-mean centred challenge and skill: flow is both above the person's own mean, relaxation skill only, anxiety challenge only, and apathy neither. Entries are differences on the raw 1 to 7 response scale from multilevel models with a random person intercept, with $p$ in parentheses. Negative values mean less pain, less interference, less attention to pain, or a lower threat value than the reference channel. The four appraisals are single momentary items; every flow variable is a secondary, analyst-constructed measure, since the source study analysed the items separately.}
\endgroup
\end{table}
}

\section{Discussion}
This study asked whether the flow construct survives contact with the daily life of patients with
chronic pain, and whether the momentary flow experience accompanies or precedes a better experience
of pain. Three findings stand out. First, the challenge-skill condition acted additively rather
than configurationally: no functional form of the match between challenge and skill improved on a
model with the two components entered separately, and between persons better-matched activity went
with less flow experience, not more. Second, the composite lost information relative to its
constituents, because absorption and skill tracked lower pain while challenge, once the others were
held constant, tracked higher pain. Third, the flow experience accompanied a better experience of
pain at the same moment, most strongly for interference, but predicted nothing at the next prompt,
and the least painful region of the challenge-skill plane was the relaxation channel rather than
flow. Each research question is discussed in turn, followed by limitations and future directions.

\subsection{Main findings}
\subsubsection{Additive rather than configural (RQ1)}
The classic account is explicitly configural: flow is not what high challenge produces, nor what
high skill produces, but what their correspondence at a high level produces
\citep{csikszentmihalyi1990flow,massimini1988systematic}. In these data that claim failed in every
form in which it was tested. The balance term was null within persons and reversed between them; a
one-sided kink at equality, a multiplicative interaction, and a full second-order response surface
each added nothing; and the four-channel quadrant scheme improved fit only marginally and was
rejected by BIC. What remained was an ordinary additive structure in which skill weighed about
twice as much as challenge.

Two features of this result deserve emphasis. The first is that the null is informative rather than
underpowered. With just under $6{,}000$ moments the balance estimate was two orders of magnitude
smaller than the elevation effect it competed with and was pinned down to within $\pm.04$, so the
finding is that there is nothing to see, not that nothing could be seen
\citep{lakens2018equivalence}. The second is that the tests do not depend on the assumption that
makes difference scores fragile. A balance term is interpretable only if a point on the challenge
scale is commensurate with a point on the skill scale, which two single diary items cannot
guarantee; the interaction and response-surface tests keep both components in their own metric and
require no such assumption \citep{edwards1993polynomial,edwards1994congruence}. They agreed with
the balance test, and the descriptive picture, flat mean flow experience across the whole range of
mismatch within a band of elevation, agreed with both.

The between-person reversal is the one positive finding in this section, and it has a clinical
reading. A patient whose activities are consistently well matched in challenge and skill is
plausibly a patient who has adjusted downward to what feels safely manageable. Challenge and skill
then correspond nicely, but at a low level, and activity avoidance scores statistically as good
balance. That interpretation connects directly to the literature on pacing, avoidance, and
endurance in chronic pain \citep{andrews2012activity,vlaeyen2000fear}, and it is a pattern one
would not expect to see in the student and worker samples in which the channel model was developed
\citep{carli1988quality}. It also suggests that the elevation term, which indexes the level of
engagement, is the one that actually varies meaningfully in this population.

These results sit alongside the wider flow literature rather than against it. In meta-analysis the
balance between challenge and skill relates to flow only modestly \citep{fong2015challenge}, and
experimental work has found the balance effect to be conditional rather than general
\citep{engeser2008flow}. The present contribution is to test the claim on its own home ground,
signal-contingent experience sampling, in a clinical sample and in several functional forms at
once.

\subsubsection{The composite adds a label where the components carry the information (RQ2)}
\citet{vanderweele2022constructed} distinguishes three worlds behind an identical composite
coefficient, and these data fall into the least comfortable of them. The composite of all four
appraisals averages over constituents that pull in opposite directions with respect to pain:
absorption and skill favourable, enjoyment redundant once the other two are present, and challenge
unfavourable. Reporting the composite alone would have hidden both the driving constituent and the
counteracting one, and would have obscured the fact that the most clinically actionable component,
challenge, is the one that carries a cost.

That challenge carries a cost is substantively plausible rather than a statistical curiosity. A
demanding moment takes something out of a person with chronic pain even when they are equal to it,
and in this sample challenge co-occurred with physical exertion. The suppression pattern is exactly
what should be expected when a variable with an unfavourable direct relation is positively
correlated with variables that have favourable ones \citep{mackinnon2000equivalence}, and it
converges with the channel finding: if challenge is unfavourable and skill favourable, the best
corner of the plane must be low-challenge and high-skill, which is what the surfaces show.

The network substitution makes the same point on a different axis. The decomposition says the
composite is carried by two of its four parts; the network says that in a dynamic system the
composite behaves indistinguishably from its dominant constituent, reproducing 22 parameters at
$r=.997$ with no change of sign or significance. Two independent lines of evidence, one
coefficient-level and one multivariate, converge on the conclusion that the flow architecture adds
a label rather than information. The recommendation that follows is concrete and transferable: when
flow is measured in experience sampling, report the constituents separately and jointly alongside
any composite, because the composite can conceal a sign reversal among them.

\subsubsection{Co-occurrence rather than analgesia (RQ3)}
The concurrent association was robust, survived adjustment for physical activation, and was ordered
as an attentional account predicts, with interference moving more than intensity
\citep{eccleston1999pain,vandamme2010keeping}. But it did not extend in time. Neither direction
produced a lag-1 effect, while the autoregressions in the same models were large, so the null is
informative about this design rather than uninterpretable. Absent temporal precedence, the relation
is co-occurrence and the direction is not identified: absorption may dampen pain, and low pain may
equally make absorption possible.

That is weaker evidence than the flow-induced-analgesia proposal requires, and it is an honest
result rather than a disappointing one. The induced-flow paradigms that motivate the proposal
measure pain during an induction \citep{gromakovskis2025flow}; the present design asks whether
naturally occurring absorption still shows a footprint an hour and a half later, and it does not.
The two findings are compatible, and the gap between them is exactly the ecological question worth
posing.

The channel result is the most directly contrary to the flow account, and also the most
clinically useful. If the goal is a less painful moment in daily life, these data point toward being
competent at something undemanding rather than toward seeking out optimal challenge. Relaxation beat
flow on interference, attention, and threat, and matched it on intensity; anxiety and apathy were
worse than both. This is a moderate, empirically grounded counterweight to the enthusiasm for
flow-based intervention: not that flow does nothing, but that in everyday life it is the skill
component that does the work, the challenge component carries a cost, and the most favourable
moments are not the flow moments.

Finally, the group average was not universal. In 4 of 63 participants the concurrent flow-pain
slope was reliably positive, and no baseline profile predicted where a given patient fell.
Baseline characteristics indexed how much flow a patient typically experienced but not how that
experience related to their pain, a dissociation between level and coupling that is an argument for
measuring the dynamics directly rather than inferring them from questionnaires
\citep{wright2020personalized,fisher2018lack,molenaar2004manifesto}.

\subsection{Limitations}
Several limitations qualify these conclusions. They concern, in turn, the measurement of the
construct, the causal status of the associations, the temporal resolution, the confounding of
challenge with exertion, the statistical power for person-level analyses, and the generality of the
sample.

First, and most importantly, the flow measure rests on four single momentary items and covers three
of the nine dimensions Csikszentmihalyi describes, leaving clear goals, unambiguous feedback, sense
of control, loss of self-consciousness, and time transformation unmeasured. It is a core flow proxy
rather than a flow scale \citep{jackson2002assessing}, and it is a secondary,
analyst-constructed measure, since the source study analysed these items separately
\citep{viane2004acceptance}. Single-item measurement also raises a specific technical objection to
the decomposition: measurement error in correlated predictors can produce sign reversals under
joint entry as an artefact. Three considerations weigh against that reading here. The reversal is
not marginal ($p=.004$ in nearly $6{,}000$ moments) and holds on all four pain outcomes; it is
theoretically interpretable, since demanding activity plausibly costs something; and it converges
with the independent channel finding, which does not depend on joint entry at all. The objection
nevertheless remains open, and replication with multi-item momentary flow scales is the way to
close it.

Second, the analysis is observational. Nothing was manipulated, and flow, pain, and activity can
influence one another through unmeasured time-varying causes such as mood, medication, or the
demands of the moment, so all estimates are within-person predictive associations rather than
causal effects \citep{rohrer2018thinking}. The channel comparison in particular contrasts moments
that participants selected rather than conditions to which they were assigned.

Third, the temporal resolution is coarse and unequal. Adjacent prompts are separated by variable
hours, no lag is carried overnight, and momentary reports are retrospective over a short window.
Flow is by definition a state that fluctuates from moment to moment, and an episode can begin and
end entirely within one interval. The absence of lagged effects is therefore open to two readings
that these data cannot separate: a measurement limitation, in which the state fluctuates faster
than the sampling grid, or a substantive null, in which flow accompanies less pain only while it
lasts. The first is the more likely and is how the result is framed here, but the second is
compatible with the induction studies, which also measure pain during rather than after the state
\citep{hamaker2017no}.

Fourth, the confounding of challenge with physical exertion could be adjusted for but not
dissolved. Challenge correlated $r=.57$ within persons with self-reported physical activation, and
although every model was estimated with activation as a covariate and none of the focal estimates
moved, the open-text activity descriptions were not coded, so demand and exertion could not be
separated substantively. This limits how far the channel ordering can be interpreted as a statement
about psychological demand rather than about physical load.

Fifth, the person-level analyses rest on 68 participants, and 63 of them contributed enough moments
for an individual slope. That is adequate for describing heterogeneity but modest for explaining
it, so the absence of trait moderation should be read as a failure to detect rather than as
evidence of no moderation.

Sixth, all data come from a single chronic-pain cohort recruited at one pain clinic, with a
predominance of women and of back and widespread pain. The channel model was developed in
non-clinical samples, and the between-person balance reversal reported here may well be specific to
a population that has reason to adjust its activities downward. Replication in other pain
conditions, and in non-clinical samples using the same nested tests, is needed before either the
structural or the substantive conclusions are generalized.

\subsection{Future directions}
Four directions follow. First, the structural tests used here cost little and should be routine.
Any study that reports a momentary flow composite can also report its constituents alone and
jointly, and any study that reports a balance term can also report the nested comparison against
an additive model and a response surface. Whether the failure of configurality found here is
specific to chronic pain or general is an empirical question that existing datasets could answer
immediately.

Second, the temporal question needs a design built for it. Burst sampling, event-contingent prompts
triggered by the start and end of activities, or continuous passive sensing of activity would
shorten the interval enough to distinguish a genuinely concurrent-only relation from one that the
ninety-minute grid cannot resolve. Only such a design can move the flow-induced-analgesia question
from co-occurrence toward precedence.

Third, the channel result invites an experimental test. Because relaxation, not flow, was the least
painful region of the plane, a within-person experiment that assigns competent but undemanding
activity against equally absorbing but demanding activity would separate the skill and challenge
components that these observational data can only decompose statistically. That is also the design
that would test whether the cost of challenge is psychological demand or physical exertion.

Fourth, the dissociation between level and coupling deserves following up. Baseline profiles
predicted how much flow a patient experienced but not how that experience related to their pain,
and four patients showed a reliably positive coupling. Whether person-specific coupling estimates
carry information beyond questionnaires, and whether they can be used to decide for whom an
activity-based intervention is appropriate, is an open and clinically consequential question
\citep{wright2020personalized}.

\section{Conclusion}
Measured moment by moment in the daily life of patients with chronic pain, the flow construct did
not hold together in the way its standard operationalization assumes: the match between challenge
and skill added nothing to the two components entered separately, in any functional form tested,
and between persons better-matched activity accompanied less flow rather than more, while the
four-item composite lost information relative to its constituents because absorption and skill
tracked lower pain and challenge, once the others were held constant, tracked higher pain. The
flow experience did accompany a better experience of pain at the same prompt, most strongly its
interference, but it predicted nothing at the next prompt while autoregressive effects in the same
models remained large, and the least painful region of the challenge-skill plane was the
low-challenge, high-skill relaxation channel rather than flow. Taken together, these results
recommend that experience-sampling research report the constituents of flow separately and jointly
rather than a composite, and they suggest that what accompanies relief in everyday life with
chronic pain is competence without demand rather than optimal challenge.

\clearpage
\bibliographystyle{apalike}
\bibliography{references}

\clearpage
\renewcommand{\thefigure}{S\arabic{figure}}
\renewcommand{\thetable}{S\arabic{table}}
\setcounter{figure}{0}
\setcounter{table}{0}

\section*{Appendix}
\addcontentsline{toc}{section}{Appendix}
\subsection*{Supplementary Figures}
\addcontentsline{toc}{subsection}{Supplementary Figures}

\begin{figure}[H]
\raggedright
\caption{The flow analytic frame.}
\label{sfig:frame}
\noindent\includegraphics[width=0.98\textwidth]{\asset{figures/supplementary/SUP_01_analytic_frame.png}}
\figurenote{Panel A shows the number of moments per participant with all four appraisals answered.
Panel B shows the item means with their standard deviations against the scale midpoint. Panel C
shows the between-person variance share of each item. Panel D compares the moments dropped by
listwise deletion with those retained; coloured bars are significant at $p<.05$. Panel E shows the
item means by prompt position within the day, restricted to positions with at least 150 moments.
Panel F shows the flow experience and pain intensity across the 14 study days.}
\end{figure}

\begin{figure}[H]
\raggedright
\caption{How the standardization rule decides every descriptive claim.}
\label{sfig:standardization}
\noindent\includegraphics[width=0.98\textwidth]{\asset{figures/supplementary/SUP_02_standardization.png}}
\figurenote{Panel A shows the median participant's share of flow moments under each rule, with the
full person range. Panel B shows the channel composition of all moments under each rule. Panel C
shows the balance and elevation effects estimated on each metric. Panel D shows the ordered
per-person prevalence under the two absolute rules. Panel E shows how often two rules agree on
whether a given moment is a flow moment. Panel F shows how many participants never reach flow under
each rule; only the absolute rules can represent absence.}
\end{figure}

\begin{figure}[H]
\raggedright
\caption{Additive-model residuals and the collinearity of the balance term.}
\label{sfig:baldiag}
\noindent\includegraphics[width=0.98\textwidth]{\asset{figures/supplementary/SUP_03_balance_diagnostics.png}}
\figurenote{Panel A shows the mean residual of the additive challenge-plus-skill model over the
challenge-skill plane; a configural effect would appear as a ridge along the dashed diagonal.
Panel B shows the coefficients of the full second-order response surface. Panel C shows the simple
slopes of balance at low, average, and high elevation inside the balance-and-elevation
parameterization. Panel D repeats the mismatch-within-elevation description on the person-centred
metric. Panel E places every nested test on a common significance scale, separating the tests of
configurality from the two diagnostics of the parameterization. Panel F shows the number of moments
in each cell of the plane.}
\end{figure}

\begin{figure}[H]
\raggedright
\caption{Temporal architecture of the construct.}
\label{sfig:temporal}
\noindent\includegraphics[width=0.98\textwidth]{\asset{figures/supplementary/SUP_04_temporal_architecture.png}}
\figurenote{Panel A shows the condition terms predicting the flow experience at the same prompt and
one prompt later, estimated on the same rows. Panel B compares the same-prompt and lagged elevation
effects directly. Panel C shows the between-person standard deviation of each lag-1 effect. Panel D
places the lag-1 flow effects next to the outcome's own autoregression from the same model. Panels
E and F show the temporal and contemporaneous networks estimated with the flow composite in place
of the absorption node; asterisks mark $p<.05$.}
\end{figure}

\begin{figure}[H]
\raggedright
\caption{Flow proneness, the flow-pain coupling, and where participants sit on the challenge-skill plane.}
\label{sfig:personlevel}
\noindent\includegraphics[width=0.98\textwidth]{\asset{figures/supplementary/SUP_05_person_level.png}}
\figurenote{Panel A shows the partial correlation of each person-level flow variable with each
baseline profile, controlling the participant's mean momentary pain; asterisks mark $p<.05$. Panel
B shows the variance in each mean momentary state added by flow proneness over the three profiles.
Panel C compares participants whose flow-pain slope is negative with the rest. Panels D and E plot
the person-specific slope against the participant's mean flow experience and against the number of
moments they contributed. Panel F places each participant on the challenge-skill plane, coloured by
their share of flow moments.}
\end{figure}

\begin{figure}[H]
\raggedright
\caption{Covariates, channel profiles, and the spread of the flow slope.}
\label{sfig:robust}
\noindent\includegraphics[width=0.98\textwidth]{\asset{figures/supplementary/SUP_06_robustness.png}}
\figurenote{Panel A shows the flow experience, balance, and elevation estimated together in the
same concurrent model. Panel B shows the effect of momentary physical activation in those models.
Panel C shows how far each focal estimate moves when day and prompt position are added as
covariates. Panel D shows the raw momentary profile of the four channels. Panel E compares the
between-person standard deviation of the flow slope with the absolute size of the fixed effect.
Panel F shows the person-centred distributions of the two focal variables.}
\end{figure}

\clearpage
\subsection*{Supplementary Tables}
\addcontentsline{toc}{subsection}{Supplementary Tables}

\begin{table}[!htbp]
\raggedright
\caption{Correlations among the four appraisals and the composite.}
\label{tab:sup-itemcorr}
\footnotesize
\begingroup
\setlength{\tabcolsep}{4pt}
\renewcommand{\arraystretch}{1.08}
\noindent
\begin{tabular*}{\textwidth}{@{\extracolsep{\fill}}lccccc@{}}
\toprule
Item & Challenge & Skill & Absorption & Enjoyment & Flow experience \\
\midrule
Challenge &  & 0.14 & 0.64 & 0.21 & 0.53 \\
Skill & 0.17 &  & 0.41 & 0.68 & 0.62 \\
Absorption & 0.53 & 0.42 &  & 0.42 & 0.88 \\
Enjoyment & 0.23 & 0.59 & 0.52 &  & 0.81 \\
Flow experience & 0.45 & 0.57 & 0.90 & 0.85 &  \\
\bottomrule
\end{tabular*}
\par\vspace{3pt}\noindent\parbox{\textwidth}{\footnotesize\textit{Note.} Below the diagonal: within-person correlations after person-mean centring. Above the diagonal: between-person correlations of the person means. The flow experience is the mean of absorption and enjoyment, so its correlations with those two items are inflated by construction.}
\endgroup
\end{table}
}
\begin{table}[!htbp]
\raggedright
\caption{Prevalence of flow moments under each classification rule.}
\label{tab:sup-prevalence}
\footnotesize
\begingroup
\setlength{\tabcolsep}{4pt}
\renewcommand{\arraystretch}{1.08}
\noindent
\begin{tabular*}{\textwidth}{@{\extracolsep{\fill}}lccccc@{}}
\toprule
Rule & Condition met & Gated flow & Median person & Person range & At zero \\
\midrule
A. Raw, both items $\geq$ 4 & 39.6 & 36.8 & 35.0 & 1.0 to 89.4 & 0/68 \\
A. Raw, both items $\geq$ 5 & 22.0 & 19.3 & 16.4 & 0.0 to 66.7 & 4/68 \\
B. Grand-mean $z$ & 30.6 & 27.8 & 26.7 & 0.0 to 75.0 & 2/68 \\
C. Within-person $z$ & 28.9 & 24.6 & 24.5 & 1.0 to 53.2 & 0/68 \\
D. Person-mean centred & 28.9 & 24.6 & 24.5 & 1.0 to 53.2 & 0/68 \\
\bottomrule
\end{tabular*}
\par\vspace{3pt}\noindent\parbox{\textwidth}{\footnotesize\textit{Note.} All shares are percentages of moments, except the last column, which counts participants who never reach flow. The gated criterion requires both the challenge-skill condition and a flow experience above the same cut. Relative rules force a within-person distribution and therefore assign above-average moments to every participant, which is why only the absolute rules can represent the absence of flow and why every prevalence claim in the text is made on the absolute metric.}
\endgroup
\end{table}
}
\begin{table}[!htbp]
\raggedright
\caption{Condition terms predicting the flow experience on every standardization metric.}
\label{tab:sup-standardization}
\footnotesize
\begingroup
\setlength{\tabcolsep}{4pt}
\renewcommand{\arraystretch}{1.08}
\noindent
\begin{tabular*}{\textwidth}{@{\extracolsep{\fill}}llcccc@{}}
\toprule
Metric & Term & $b$ & SE & 95\% CI & $p$ \\
\midrule
D. Person-mean centered (primary) & Balance & 0.004 & 0.018 & [-0.031, 0.039] & .826 \\
D. Person-mean centered (primary) & Elevation & 0.743 & 0.023 & [0.698, 0.788] & <.001 \\
C. Within-person z & Balance & 0.005 & 0.021 & [-0.036, 0.046] & .806 \\
C. Within-person z & Elevation & 0.857 & 0.021 & [0.816, 0.898] & <.001 \\
B. Grand-mean z & Balance & -0.025 & 0.025 & [-0.074, 0.024] & .315 \\
B. Grand-mean z & Elevation & 0.890 & 0.022 & [0.847, 0.933] & <.001 \\
A. Raw absolute scale & Balance & -0.090 & 0.021 & [-0.131, -0.049] & <.001 \\
A. Raw absolute scale & Elevation & 0.768 & 0.020 & [0.729, 0.807] & <.001 \\
\bottomrule
\end{tabular*}
\par\vspace{3pt}\noindent\parbox{\textwidth}{\footnotesize\textit{Note.} Person-mean centring is the primary metric. Across all four metrics the elevation effect stays between .74 and .89 and the balance effect is never positive.}
\endgroup
\end{table}
}
\begin{table}[!htbp]
\raggedright
\caption{Fixed effects of every nested condition model.}
\label{tab:sup-condmodels}
\footnotesize
\begingroup
\setlength{\tabcolsep}{4pt}
\renewcommand{\arraystretch}{1.08}
\noindent
\begin{tabular*}{\textwidth}{@{\extracolsep{\fill}}llcccc@{}}
\toprule
Model & Term & $b$ & SE & 95\% CI & $p$ \\
\midrule
challenge + skill & Challenge & 0.285 & 0.008 & [0.270, 0.300] & <.001 \\
challenge + skill & Skill & 0.507 & 0.010 & [0.488, 0.526] & <.001 \\
challenge + skill + balance & Challenge & 0.284 & 0.008 & [0.269, 0.300] & <.001 \\
challenge + skill + balance & Skill & 0.508 & 0.010 & [0.488, 0.527] & <.001 \\
challenge + skill + balance & Balance & -0.009 & 0.011 & [-0.031, 0.014] & .452 \\
challenge + skill + overload & Challenge & 0.276 & 0.014 & [0.248, 0.304] & <.001 \\
challenge + skill + overload & Skill & 0.516 & 0.016 & [0.485, 0.548] & <.001 \\
challenge + skill + overload & Overload kink & 0.017 & 0.023 & [-0.028, 0.062] & .452 \\
challenge x skill & Challenge & 0.285 & 0.008 & [0.270, 0.300] & <.001 \\
challenge x skill & Skill & 0.507 & 0.010 & [0.488, 0.526] & <.001 \\
challenge x skill & Challenge $\times$ skill & 0.000 & 0.005 & [-0.010, 0.010] & .970 \\
quadratic response surface & Challenge & 0.289 & 0.008 & [0.273, 0.304] & <.001 \\
quadratic response surface & Skill & 0.508 & 0.011 & [0.487, 0.530] & <.001 \\
quadratic response surface & Challenge$^2$ & -0.008 & 0.004 & [-0.016, 0.000] & .051 \\
quadratic response surface & Skill$^2$ & 0.003 & 0.005 & [-0.008, 0.013] & .622 \\
quadratic response surface & Challenge $\times$ skill & 0.001 & 0.005 & [-0.009, 0.012] & .804 \\
quadratic surface + balance & Challenge & 0.295 & 0.008 & [0.279, 0.311] & <.001 \\
quadratic surface + balance & Skill & 0.503 & 0.011 & [0.482, 0.525] & <.001 \\
quadratic surface + balance & Challenge$^2$ & -0.034 & 0.007 & [-0.048, -0.020] & <.001 \\
quadratic surface + balance & Skill$^2$ & -0.019 & 0.007 & [-0.033, -0.005] & .010 \\
quadratic surface + balance & Balance & -0.130 & 0.030 & [-0.188, -0.071] & <.001 \\
quadratic surface + balance & Challenge $\times$ skill & 0.048 & 0.012 & [0.024, 0.071] & <.001 \\
balance x elevation & Balance & 0.020 & 0.012 & [-0.003, 0.043] & .084 \\
balance x elevation & Elevation & 0.810 & 0.019 & [0.773, 0.847] & <.001 \\
balance x elevation & Balance $\times$ elevation & 0.043 & 0.010 & [0.023, 0.063] & <.001 \\
\bottomrule
\end{tabular*}
\par\vspace{3pt}\noindent\parbox{\textwidth}{\footnotesize\textit{Note.} All models predict the flow experience from person-mean centred predictors with a random person intercept, estimated by maximum likelihood. The ``quadratic surface + balance'' model is reported for completeness; its balance coefficient is not interpretable in isolation because balance correlates .93 with the squared discrepancy that the surface already contains.}
\endgroup
\end{table}
}
\begin{table}[!htbp]
\raggedright
\caption{Diagnostics of the balance parameterization.}
\label{tab:sup-baldiag}
\footnotesize
\begingroup
\setlength{\tabcolsep}{4pt}
\renewcommand{\arraystretch}{1.08}
\noindent\textit{Panel A.} Why the balance term is not interpretable on its own.\par\vspace{2pt}
\begin{tabular*}{\textwidth}{@{\extracolsep{\fill}}lc@{}}
\toprule
Quantity & Value \\
\midrule
$r$(balance, $-(C-S)^2$) & 0.927 \\
$r$(balance, $C^2$) & -0.506 \\
$r$(balance, $S^2$) & -0.324 \\
$r$(balance, elevation) & 0.030 \\
$b$ challenge (additive model) & 0.285 \\
$b$ skill (additive model) & 0.507 \\
\bottomrule
\end{tabular*}
\par\vspace{6pt}\noindent\textit{Panel B.} Simple slopes of balance across elevation.\par\vspace{2pt}
\begin{tabular*}{\textwidth}{@{\extracolsep{\fill}}lccccc@{}}
\toprule
Elevation & Value (centred) & Balance $b$ & SE & 95\% CI & $p$ \\
\midrule
-1 SD & -1.179 & -0.031 & 0.016 & [-0.063, 0.002] & .064 \\
person mean & 0.000 & 0.020 & 0.012 & [-0.003, 0.043] & .084 \\
+1 SD & 1.179 & 0.070 & 0.017 & [0.037, 0.103] & <.001 \\
\bottomrule
\end{tabular*}
\par\vspace{6pt}\noindent\textit{Panel C.} The two comparisons in which balance appears to matter.\par\vspace{2pt}
\begin{tabular*}{\textwidth}{@{\extracolsep{\fill}}lcccc@{}}
\toprule
Comparison & $\chi^2$ & df & $p$ & $\Delta$ AIC \\
\midrule
balance added to the quadratic surface & 18.87 & 1 & <.001 & 16.9 \\
balance x elevation inside the balance-elevation form & 17.60 & 1 & <.001 & 15.6 \\
\bottomrule
\end{tabular*}
\par\vspace{3pt}\noindent\parbox{\textwidth}{\footnotesize\textit{Note.} $C$ is challenge and $S$ skill, both person-mean centred. Balance is a monotone transformation of the absolute discrepancy and therefore nearly collinear with the squared discrepancy that a quadratic surface already contains; the balance-and-elevation form additionally constrains challenge and skill to equal weight, which the free additive estimates reject. Both comparisons in Panel C are properties of those parameterizations rather than evidence of configurality, which is why no configural form beats the additive model in \Cref{tab:condition}.}
\endgroup
\end{table}
}
\begin{table}[!htbp]
\raggedright
\caption{Does the challenge-skill condition precede the flow experience?}
\label{tab:sup-condlag}
\footnotesize
\begingroup
\setlength{\tabcolsep}{4pt}
\renewcommand{\arraystretch}{1.08}
\noindent
\begin{tabular*}{\textwidth}{@{\extracolsep{\fill}}llcccc@{}}
\toprule
Model & Term & $b$ & SE & $p$ & $n$ obs \\
\midrule
Lag 1 & Balance at $t-1$ & -0.007 & 0.017 & .696 & 5,065 \\
Lag 1 & Elevation at $t-1$ & 0.052 & 0.028 & .062 & 5,065 \\
Lag 1 & Flow experience at $t-1$ & 0.108 & 0.019 & <.001 & 5,065 \\
Same beep & Balance, same beep & -0.001 & 0.022 & .963 & 5,065 \\
Same beep & Elevation, same beep & 0.856 & 0.022 & <.001 & 5,065 \\
\bottomrule
\end{tabular*}
\par\vspace{3pt}\noindent\parbox{\textwidth}{\footnotesize\textit{Note.} Both models are fitted on the same rows so the lagged and same-beep estimates are directly comparable. The lag-1 model controls the flow experience's own autoregression; predictors are person-standardized and lags are built within day.}
\endgroup
\end{table}
}
\begin{table}[!htbp]
\raggedright
\caption{Substituting the flow composite for the single absorption node.}
\label{tab:sup-network}
\footnotesize
\begingroup
\setlength{\tabcolsep}{4pt}
\renewcommand{\arraystretch}{1.08}
\noindent
\begin{tabular*}{\textwidth}{@{\extracolsep{\fill}}lccccc@{}}
\toprule
Edge & Benchmark (absorption) & $p$ & With the composite & $p$  & Difference \\
\midrule
Pain(t-1) $\rightarrow$ activity node(t) & -0.048 & .065 & -0.021 & .413 & 0.027 \\
Activity node(t-1) $\rightarrow$ Pain(t) & 0.003 & .809 & 0.000 & .991 & -0.003 \\
Pain - activity node (contemporaneous) & -0.064 & .003 & -0.066 & .002 & -0.002 \\
Attention - activity node (contemporaneous) & -0.022 & .188 & -0.027 & .109 & -0.005 \\
\bottomrule
\end{tabular*}
\par\vspace{3pt}\noindent\parbox{\textwidth}{\footnotesize\textit{Note.} Both networks are multilevel vector autoregressive models on four nodes: momentary pain, its threat value, attention to pain, and either goal absorption (benchmark) or the flow composite. Only the edges that involve the substituted node can change; none changes sign or significance.}
\endgroup
\end{table}
}
\begin{table}[!htbp]
\raggedright
\caption{Momentary profile of the four challenge-skill channels.}
\label{tab:sup-channelprofiles}
\footnotesize
\begingroup
\setlength{\tabcolsep}{4pt}
\renewcommand{\arraystretch}{1.08}
\noindent
\begin{tabular*}{\textwidth}{@{\extracolsep{\fill}}lcccccccc@{}}
\toprule
Channel & $n$ & \% & Pain & Interference & Attention & Threat & Flow exp. & Activation \\
\midrule
Flow & 1,314 & 22.0 & 4.90 & 4.15 & 2.99 & 2.34 & 5.79 & 4.33 \\
Relaxation & 2,214 & 37.0 & 4.29 & 3.52 & 2.69 & 2.09 & 4.64 & 2.47 \\
Anxiety & 516 & 8.6 & 4.83 & 4.28 & 3.11 & 2.23 & 4.44 & 4.44 \\
Apathy & 1,932 & 32.3 & 4.49 & 3.94 & 3.00 & 2.31 & 3.33 & 2.46 \\
\bottomrule
\end{tabular*}
\par\vspace{3pt}\noindent\parbox{\textwidth}{\footnotesize\textit{Note.} Unadjusted momentary means on the raw 1 to 7 scale under the strict absolute rule (both condition items at least 5). Activation is the self-reported physical activation item, which is far higher in the two high-challenge channels and is the reason every model reported in the text is also estimated with activation as a covariate.}
\endgroup
\end{table}
}
\begin{table}[!htbp]
\raggedright
\caption{Person-level flow proneness and the baseline profiles.}
\label{tab:sup-personlevel}
\footnotesize
\begingroup
\setlength{\tabcolsep}{4pt}
\renewcommand{\arraystretch}{1.08}
\noindent\textit{Panel A.} Person-level flow variables against the baseline profiles.\par\vspace{2pt}
\begin{tabular*}{\textwidth}{@{\extracolsep{\fill}}llccccc@{}}
\toprule
Person-level flow variable & Baseline profile & $n$ & $r$ & $p$ & Partial $r$ & $p$  \\
\midrule
Flow proneness (strict) & Threat & 67 & -0.092 & .459 & -0.146 & .238 \\
Flow proneness (strict) & Biomedical & 67 & 0.294 & .016 & 0.048 & .702 \\
Flow proneness (strict) & Personality & 64 & 0.045 & .722 & -0.097 & .444 \\
Mean flow experience & Threat & 67 & -0.268 & .028 & -0.274 & .025 \\
Mean flow experience & Biomedical & 67 & 0.032 & .798 & 0.003 & .982 \\
Mean flow experience & Personality & 64 & -0.264 & .035 & -0.285 & .022 \\
Mean challenge & Threat & 67 & -0.047 & .707 & -0.096 & .440 \\
Mean challenge & Biomedical & 67 & 0.166 & .178 & -0.131 & .290 \\
Mean challenge & Personality & 64 & 0.157 & .216 & 0.030 & .812 \\
Mean skill & Threat & 67 & -0.114 & .360 & -0.117 & .344 \\
Mean skill & Biomedical & 67 & 0.117 & .345 & 0.113 & .361 \\
Mean skill & Personality & 64 & -0.172 & .174 & -0.185 & .142 \\
Flow-pain coupling & Threat & 62 & -0.119 & .355 & -0.171 & .185 \\
Flow-pain coupling & Biomedical & 62 & 0.138 & .284 & -0.094 & .466 \\
Flow-pain coupling & Personality & 59 & -0.016 & .905 & -0.141 & .287 \\
\bottomrule
\end{tabular*}
\par\vspace{6pt}\noindent\textit{Panel B.} Variance in the mean momentary states added by flow proneness.\par\vspace{2pt}
\begin{tabular*}{\textwidth}{@{\extracolsep{\fill}}lcccc@{}}
\toprule
Mean momentary state & $n$ & $R^2$ profiles & $R^2$ + flow proneness & $\Delta R^2$ \\
\midrule
Attention to pain & 64 & 0.196 & 0.229 & 0.033 \\
Pain intensity & 64 & 0.362 & 0.443 & 0.081 \\
Pain interference & 64 & 0.355 & 0.388 & 0.033 \\
Threat value & 64 & 0.408 & 0.417 & 0.009 \\
\bottomrule
\end{tabular*}
\par\vspace{6pt}\noindent\textit{Panel C.} Participants whose flow-pain slope is negative against the rest.\par\vspace{2pt}
\begin{tabular*}{\textwidth}{@{\extracolsep{\fill}}lcccccc@{}}
\toprule
Person-level variable & $n$ negative & $n$ other & Mean, negative & Mean, other & Cohen's $d$ & $p$ \\
\midrule
Threat & 45 & 17 & 0.059 & -0.271 & 0.398 & .151 \\
Biomedical & 45 & 17 & -0.071 & 0.272 & -0.460 & .131 \\
Personality & 42 & 17 & -0.002 & -0.133 & 0.129 & .654 \\
Flow proneness (strict) & 45 & 17 & 0.161 & 0.293 & -0.739 & .027 \\
Mean pain & 45 & 17 & 4.313 & 5.073 & -0.690 & .029 \\
\bottomrule
\end{tabular*}
\par\vspace{3pt}\noindent\parbox{\textwidth}{\footnotesize\textit{Note.} Partial correlations control the participant's mean momentary pain. The threat profile combines the Pain Catastrophizing Scale and the Pain Vigilance and Awareness Questionnaire, the biomedical profile pain severity and duration, and the personality profile NEO-FFI neuroticism. Flow proneness is the share of a participant's moments meeting the strict absolute criterion. No profile predicts the person-specific flow-pain slope, which is the contrast the text draws with the level effects.}
\endgroup
\end{table}
}
\begin{table}[!htbp]
\raggedright
\caption{Robustness of the flow analyses.}
\label{tab:sup-robustness}
\footnotesize
\begingroup
\setlength{\tabcolsep}{4pt}
\renewcommand{\arraystretch}{1.08}
\noindent\textit{Panel A.} The flow-pain estimate at four compliance floors.\par\vspace{2pt}
\begin{tabular*}{\textwidth}{@{\extracolsep{\fill}}lccccc@{}}
\toprule
Compliance floor & $n$ persons & $n$ moments & $b$ & SE & $p$ \\
\midrule
$\geq$ 20 moments & 68 & 5,969 & -0.0937 & 0.0200 & <.001 \\
$\geq$ 30 moments & 67 & 5,944 & -0.0938 & 0.0202 & <.001 \\
$\geq$ 50 moments & 63 & 5,772 & -0.0912 & 0.0208 & <.001 \\
$\geq$ 70 moments & 61 & 5,653 & -0.0883 & 0.0212 & <.001 \\
\bottomrule
\end{tabular*}
\par\vspace{6pt}\noindent\textit{Panel B.} Day and prompt position added as level-1 covariates.\par\vspace{2pt}
\begin{tabular*}{\textwidth}{@{\extracolsep{\fill}}llccccc@{}}
\toprule
Model & Focal term & $b$ without & $b$ with & $p$ & $p$ beep & $p$ day \\
\midrule
V2: condition $\rightarrow$ experience & elevation\_w & 0.7443 & 0.7456 & <.001 & <.001 & .959 \\
Flow $\rightarrow$ pain (activity-adjusted) & FLOWEXP\_w & -0.0937 & -0.1010 & <.001 & <.001 & .079 \\
\bottomrule
\end{tabular*}
\par\vspace{6pt}\noindent\textit{Panel C.} Moments lost to listwise deletion.\par\vspace{2pt}
\begin{tabular*}{\textwidth}{@{\extracolsep{\fill}}lccccc@{}}
\toprule
Measure & Mean, retained & Mean, dropped & $n$ dropped & Cohen's $d$ & $p$ \\
\midrule
Pain intensity & 4.53 & 4.36 & 286 & -0.115 & .052 \\
Pain interference & 3.86 & 3.79 & 272 & -0.043 & .476 \\
Attention to pain & 2.89 & 2.95 & 255 & 0.034 & .593 \\
Threat value & 2.23 & 2.44 & 268 & 0.151 & .019 \\
Negative affect & 1.92 & 2.14 & 280 & 0.193 & .002 \\
Positive affect & 4.23 & 4.27 & 286 & 0.033 & .583 \\
Physical activation & 3.05 & 2.62 & 274 & -0.235 & <.001 \\
\bottomrule
\end{tabular*}
\par\vspace{3pt}\noindent\parbox{\textwidth}{\footnotesize\textit{Note.} Panel A repeats the activity-adjusted concurrent model of pain intensity with a stricter minimum number of completed moments per participant. In Panel B the focal estimate is compared with and without day and prompt position; prompt position is itself significant but moves no focal estimate. Panel C compares the moments retained in the flow analytic frame with the 4.6\% dropped for a missing appraisal.}
\endgroup
\end{table}
}

\end{document}
